
\label{ch:conclusions}


\begin{ataglance}[How to read this chapter]
This chapter closes the argument. It states what the thesis produced, answers
each aim in turn, marks the boundary of every claim and says what should be
done next. The order runs from what the work claims to what it cannot claim.

\S\ref{sec:contribution} sets out the contribution, with the axiomatic
realisation model first, because the other three contributions exist to give it
something to be tested against, and ends with the systems-engineering framing
of the instrument. \S\ref{sec:aims} answers the four aims one at a time, and
marks with placeholders those answers that wait on the analysis in
Chapters~\ref{ch:results1} and~\ref{ch:results2}. \S\ref{sec:conc-questions}
takes the research questions of Chapter~\ref{ch:introduction} one by one and
says how much of each answer the thesis already holds. \S\ref{sec:hypotheses}
does the same for the hypotheses, and does it more briefly, since a hypothesis
has fewer ways of being half-answered than an aim. \S\ref{sec:scope} is the
section that costs the thesis most: one cohort, one country, one measurement
occasion, an estimated predictor, and normative choices that no quantity of
data will settle. \S\ref{sec:conc-use} says what the instrument may be used
for now and what must wait. \S\ref{sec:future} says what is worth doing next and
in what order. \S\ref{sec:conc-systems} reads the conclusions through the
systems view of Chapter~\ref{ch:systems} and sets out the monitoring and
governance a deployed score would need. \S\ref{sec:final} states the argument in
a few paragraphs.

Two technical passages carry plain readings.
Derivation~\ref{der:bridge} explains why the conceptual formula is not the
scoring equation, and Derivation~\ref{der:uncertainty} explains how much
predictive strength is lost by treating an estimated score as though it had
been measured. Both use the reference profiles of Table~\ref{tab:profiles}.
\end{ataglance}

\begin{quote}
\small\itshape
Note to the reader. The verdicts on the empirical hypotheses depend on the analysis reported in Chapters~\ref{ch:results1} and~\ref{ch:results2}. Statements below that depend directly upon those results contain placeholders to be completed once the analysis is finalised. This note is removed in the final version.
\end{quote}

\noindent What may this thesis claim, and on what terms? The question has to be
answered before the analysis is complete. The chapter must therefore say in
advance what each possible result would mean, so that no verdict can be shaped
to fit whatever the data return; the rules for reading a null, a partial fit or
a strong association are fixed while no one knows which will be needed. There is a second task, which the systems
chapter (Chapter~\ref{ch:systems}) made unavoidable. A score that is used acts
on the people it describes. The conclusions must cover how the instrument would
be used and watched, as well as what it measures and predicts.

This thesis set out to develop, validate and evaluate a multidimensional health metric, the $H_{\text{potential}}$, that conceptualises health as future capability rather than solely as current disease status, accumulated impairment or present quality of life. It also examined whether this general framework could be translated coherently into a disease-specific context, using haemophilia as the first application. The work was motivated by a fundamental problem in contemporary health measurement: many established instruments describe important aspects of health, yet few provide an integrated account of the biological, social and structural conditions that shape an individual's capacity to maintain health over time.

The thesis addressed this problem through a sequential programme of conceptual, measurement, prognostic and translational research. It began by developing a theory of health capability, proceeded to operationalise that theory as a latent measurement model, evaluated the resulting construct against future clinical outcomes, and concluded by setting out a structured pathway for adaptation to haemophilia. This sequence was deliberate. It ensured that the construct was defined independently of the outcomes it was later asked to predict and that disease-specific translation occurred only after the general framework had been established.

The central proposition advanced by the thesis is that health should not be understood exclusively as a static state. Current disease burden, disability, symptoms and health-related quality of life all provide valuable information, but they do not necessarily reveal the same thing as an individual's remaining capacity to sustain health, adapt to biological challenge and benefit from available social and healthcare resources. Two individuals may appear similar when assessed through conventional measures of current status while differing substantially in their future health trajectories because their physiological reserve, social support and healthcare environments differ. The $H_{\text{potential}}$ was developed to provide a formal representation of this difference.

The framework therefore treats health as a latent capability arising from three domains: Biomedical Potential, the physiological and functional resources available to the individual; Social Modifiers, the interpersonal, behavioural and material conditions that enable or constrain their use; and Systemic Burdens, the structural and health-system barriers that decide whether support and treatment can be realised. The domains can be told apart in analysis but are not independent in practice; their interaction decides how far current capacities become future health.

% =====================================================================
\section{Contribution of the Thesis}
\label{sec:contribution}
% =====================================================================

\noindent The central contribution of the thesis is a formal one: the axiomatic Capability Realisation Model and the integrated, equity-sensitive policy calculation built upon it (Chapter~\ref{chap:crf}). Three further connected contributions support and surround it: a conceptual contribution to the theory of health measurement, a methodological contribution to the development and evaluation of multidimensional health metrics, and a translational contribution concerning adaptation from a general population to a disease-specific setting. Their relationship is shown in Figure~\ref{fig:contribution_triangle}, which places the conceptual, methodological and translational contributions at the corners and the formal model they serve at the centre. A fifth contribution, the systems-engineering framing of the instrument set out in Chapter~\ref{ch:systems}, lies outside the triangle because it concerns how the whole would be specified, checked and governed in use rather than what any one part claims.

The formal model is taken first because the three that follow exist to give
it something to be tested against.

\subsection{The Central Contribution: an Axiomatic Capability Realisation Model}

The formal heart of the thesis is the Capability Realisation Model developed in Chapter~\ref{chap:crf}. Prior multidimensional health measures, including the earlier heuristic $(L\cdot Q)S-B$, assert an aggregation rule. The realisation model instead states the properties a realisation measure must satisfy, chief among them non-compensation, that complete social exclusion or complete structural obstruction cannot be cancelled by biomedical advantage, and derives the admissible family from those axioms. Its primary form, $H(P,S,R)=P\,S^{\alpha}R^{1-\alpha}$ with $R=1-B$, is characterised uniquely by constant contextual elasticities, and from it several results follow as theorems rather than assertions: biomedical--context complementarity, which is the mathematical form of the intersectional claim and which an additive index cannot represent; an exact decomposition of any capability gap into biomedical, social and structural parts without a residual term; a budget-constrained policy gradient that ranks marginal, model-implied gains from biomedical, social and burden-reducing interventions; and an equity-sensitive population functional whose inequality-aversion \citep{atkinson_1970} parameter makes the normative weighting of low-capability groups explicit rather than hidden. Together these turn an individual capability estimate into an integrated, distribution-aware policy calculation. The model is falsifiable: the chapter pre-specifies the head-to-head comparison, against an additive composite and the outcome-independent factor score, that would reject or narrow it, with its parameters frozen without reference to the outcomes. This is the contribution the rest of the thesis is built to test.

\begin{reading}[Non-compensation with actual numbers]
Set $\alpha=0.5$ and take Profile~B. With structure favourable the profile is
$(0.30,0.70,0.85)$, and
$H=0.30\times\sqrt{0.70}\times\sqrt{0.85}=0.23$. With structure unfavourable
only the third coordinate changes, to $0.20$:
$H=0.30\times\sqrt{0.70}\times\sqrt{0.20}=0.11$. Biomedical potential is
unchanged, and realised capability has halved.

An equally weighted additive index records the same move as a fall from $0.62$
to $0.40$, a loss of about a third. The difference between the two arithmetics
is the whole argument. The additive form lets the domains trade: it says the
structural loss can be offset if the biomedical score rises far enough. The
multiplicative form says how far. To return to $0.23$ with structure
unfavourable, $P$ would have to rise from $0.30$ to $0.62$.
Non-compensation is not a preference about weighting. It is the claim that some
deficits cannot be purchased from another domain at any price.
\end{reading}

\begin{figure}[htbp]
\centering
\begin{tikzpicture}
\node[hpkey, text width=3.6cm] (core) at (0,0) {Common research architecture};
\node[hpbox, text width=3.9cm] (c1) at (90:3.0) {Conceptual model of health as future capability};
\node[hpbox, text width=3.9cm] (c2) at (210:3.6) {Two-stage method that separates measurement from prediction};
\node[hpbox, text width=3.9cm] (c3) at (330:3.6) {Translational route from a disease-independent metric to disease-specific use};
\foreach \c in {c1,c2,c3} \draw[hparrow] (\c) -- (core);
\draw[hpdash] (c1) -- (c2) -- (c3) -- (c1);
\end{tikzpicture}

\captionsetup{justification=justified, width=\linewidth}
\caption{The three connected contributions of the thesis.}
\label{fig:contribution_triangle}
\end{figure}

A set of axioms constrains the form a measure may take. It says nothing
about what the measure is of, and that question was settled first.

\subsection{A Conceptual Model of Health as Future Capability}

The first contribution is the development of a theoretically grounded model of health as future capability. The framework brings together ideas from the capability approach, healthy ageing, frailty science and the social determinants of health. These literatures have often developed in parallel: the capability approach emphasises the real opportunities available to individuals; healthy ageing focuses on the interaction between intrinsic capacity \citep{who_2015_ageing,cesari_2018} and environment; frailty research examines vulnerability and physiological reserve; and the social-determinants literature demonstrates how social position, material conditions and healthcare access shape health outcomes. The present thesis integrates these insights within a common measurement architecture.

This integration matters because conventional measures frequently isolate these dimensions. Frailty indices primarily capture biological and functional deficits. Quality-of-life instruments describe current experience. Comorbidity indices quantify diagnosed disease burden. Economic measures such as QALYs support comparative valuation and resource allocation. Social and structural variables are often included as external covariates after the main health measure has already been defined. The $H_{\text{potential}}$ instead treats biological, social and systemic influences as constitutive elements of a single latent health capability.

The conceptual contribution is therefore not simply the addition of more variables to an existing index. It is a change in the object being measured. The target is the underlying capability from which future health trajectories may emerge. The framework also clarifies that health capability is neither exclusively biological nor reducible to social advantage. Biomedical Potential remains fundamental because physiological and functional reserve constrain what an individual can sustain, yet biological reserve cannot be interpreted apart from the social and systemic conditions through which it is expressed. Conversely, favourable social and healthcare environments cannot fully compensate for every form of severe biological impairment. The model therefore locates future capability in the interaction among the domains, rejecting both biomedical reductionism and the view that health is wholly socially determined.

Deciding what to measure and showing that it can be measured are separate
jobs. The thesis was arranged so that the second could fail without taking the
first down with it.

\subsection{A Method for Separating Measurement from Prediction}

The second contribution is methodological. The thesis operationalises the conceptual framework as a latent construct estimated through structural equation modelling and then evaluates that construct in a separate prognostic model. This two-stage design was chosen to reduce a common form of circularity in composite measure development. Many prognostic indices are constructed by selecting variables and estimating weights according to their observed relationships with the outcome they are intended to predict. Such approaches may be entirely appropriate when the sole objective is prediction; however, they make it difficult to distinguish the measurement of a theoretical construct from the optimisation of a prediction equation, and the resulting score may predict effectively while remaining uncertain in conceptual meaning.

The present thesis separates these tasks. The measurement model is estimated without using the future clinical outcome to define the construct, and its evaluation is based upon factor structure, model fit, reliability and construct validity. Only after the factor scores are derived is the $H_{\text{potential}}$ entered into survival models predicting mortality and hospitalisation. Prognostic evaluation therefore tests whether an independently defined construct possesses future predictive value. This separation strengthens interpretability: if the construct predicts future outcomes, the result is evidence that the theoretically defined latent capability has prognostic relevance; if it does not, the measurement model does not become meaningless, but rather the result indicates that the construct, as operationalised, may not support the intended prognostic use. Measurement and prediction can therefore succeed or fail independently, allowing more precise scientific conclusions.

The design also provides a reusable methodological template. Researchers developing multidimensional measures in other fields can define the construct first, estimate it through an appropriate measurement model, and evaluate its predictive performance in a distinct analytical stage. A further methodological contribution lies in the breadth of the prognostic evaluation: the thesis does not treat a statistically significant hazard ratio as sufficient evidence of model value, but instead assesses association, discrimination, calibration, comparative model fit, incremental value, decision-analytic utility and optimism-corrected internal performance.

The template was built on one cohort in one country. Whether it transfers
is a question about a different population, and haemophilia is where the thesis
asks it.

\subsection{A Translational Route from General to Disease-Specific Use}

The third contribution is the development of a structured route from a disease-independent measure to disease-specific application. Haemophilia was selected as the first translational context because improvements in treatment have transformed it from a condition dominated by early mortality into a lifelong chronic disease in which ageing, musculoskeletal function, multimorbidity, social support, treatment access and healthcare-system stability all influence long-term outcomes.

The haemophilia application preserves the three latent domains while replacing general indicators with disease-relevant measures. Joint health, bleeding phenotype, functional mobility and inhibitor status represent Biomedical Potential. Access to comprehensive care, health literacy, caregiver support and self-management represent Social Modifiers. Treatment availability, supply-chain stability, financial coverage and structural barriers represent Systemic Burdens. This is not presented as a completed haemophilia instrument. No haemophilia data were analysed in the thesis, and no disease-specific factor loadings or prognostic parameters were estimated. The contribution instead lies in the translation method: indicator mapping, expert consultation through a modified Delphi process, re-estimation of the measurement model, assessment of measurement invariance \citep{meredith_1993,vandenberg_lance_2000}, recalibration and external prognostic validation. The domain structure is the element proposed for preservation; the empirical parameters must be re-estimated in the target population.

The three contributions so far describe the instrument at rest: what it means,
how it is estimated, how it moves to a new population. None says what happens
once a score is computed and someone acts on it. That is the subject of the
next contribution.

\subsection{An Engineered Instrument: Requirements, Failure Modes and Control}
\label{sec:conc-engineered}

A measure that will be used to sort people is a built thing, and it can fail in
the ways built things fail. Chapter~\ref{ch:systems} treats the
$H_{\text{potential}}$ in that way, and the treatment has three parts.

The first is a set of requirements traced to tests
(\S\ref{sec:sys-pipeline}). Each property the instrument must have is written as
a requirement that can be checked: the construct is defined without the
outcomes; the realisation function satisfies its axioms; $\alpha$, $\rho$ and
the normalisation bounds are frozen before the outcomes are seen; every stage is
refitted inside each bootstrap replicate; reported performance includes
calibration as well as discrimination. Each requirement points to what checks
it, either the verification code of Appendix~\ref{app:verification}, which
asserts every computable statement and worked example in the text, or a step in
the pre-specified analysis plan. Tracing requirements to their verification in
this way is standard practice in systems engineering \citep{incose_2015}. Its
value here is that a requirement with no test is a hope, and the record shows
which requirements are which. Appendix~\ref{app:verification} is explicit that
its suite checks mathematics and arithmetic, not medicine; the requirements that
need data remain open until the analysis is run.

The second part is a record of failure modes and their effects
(\S\ref{sec:sys-fmea}), in the form used in engineering practice
\citep{stamatis_2003}. Each way the pipeline could return a wrong answer, from a
miscoded indicator to a factor score too unreliable to carry a null, is listed
with its effect on the conclusions and the control that guards against it. The
worked example of a null at low reliability later in this chapter is one entry
from that record, worked through in full. Appendix~\ref{app:datachecks} is
another: it lists the changes the data checks forced on the build, two of which
would have changed the results had they gone unnoticed.

The third part is a control structure (\S\ref{sec:sys-control}). It names who
would act on the score, what they would do, and what feedback would tell them
that an action had gone wrong \citep{leveson_2011}. No psychometric or
prognostic test supplies this. A score can pass every statistical test and still
cause harm if those who act on it get no signal when it misleads them.

This framing adds no finding. It states, before any finding, what must be true
for a finding to be safe to act on.

Each of these contributions is modest taken alone. What they do together is
the reason for the arrangement.

\subsection{The Combined Contribution}

Taken together, these contributions place a falsifiable formal model at the centre and build the empirical and translational apparatus needed to test it. The axiomatic realisation model specifies, from first principles, how biomedical potential, social enablement and structural freedom combine into realised capability and into the policy calculations that follow. The conceptual model specifies what the construct is intended to mean. The measurement model tests whether that construct can be represented empirically and supplies the outcome-independent inputs and comparator for the realisation model. The prognostic model tests whether the resulting scores relate to future clinical outcomes and carries the head-to-head comparison that can reject or narrow the realisation model. The translational framework tests whether the same architecture can be adapted to a new clinical setting without being rebuilt from first principles. The systems framing states what must surround the instrument, in tests, records and feedback, before any of these results is acted on. The arrangement also lets each part fail without bringing down the parts before it. The combined contribution is therefore larger than the production of a single score: it is an axiomatic, distribution-aware model of health capability, and a research architecture for subjecting such a model to empirical and translational scrutiny.

That is what the thesis set out to build. What it undertook to demonstrate was
written down in advance as four aims, and each can now be answered.

% =====================================================================
\section{Answering the Research Aims}
\label{sec:aims}
% =====================================================================

\noindent The thesis addressed four research aims, each corresponding to a distinct stage in the development of the $H_{\text{potential}}$. In brief:

\begin{enumerate}
  \item \textbf{Develop the conceptual framework.} Chapter~\ref{ch:theory} set out the three-domain model and the conceptual formula that motivates it, and the bridge from that formula to the latent construct.
  \item \textbf{Develop and validate the measurement model.} Chapter~\ref{ch:results1} estimated the second-order structural equation model and evaluated its fit, reliability, and construct validity.
  \item \textbf{Evaluate prognostic performance.} Chapter~\ref{ch:results2} tested the $H_{\text{potential}}$ against five-year mortality and first overnight hospital stay and compared it with established measures.
  \item \textbf{Assess translational relevance to haemophilia.} Chapter~\ref{ch:translational} mapped the domains to haemophilia indicators, set out the Delphi \citep{hasson_2000} step for face validity, and gave a roadmap for external validation.
\end{enumerate}

\noindent The remainder of this section records the verdict on each aim in turn. The
first aim is conceptual and can be answered in full here; the three that follow
depend in part on results reported elsewhere.

\subsection{Aim 1: Develop the Conceptual Framework}

The first aim was to develop a theoretically coherent framework for measuring health as future capability. This aim was addressed in Chapter~\ref{ch:theory}, which established the conceptual basis of the $H_{\text{potential}}$ and defined its three domains. The chapter argued that current disease status alone is insufficient to represent future health capability, and showed that the capability approach \citep{sen_1993,venkatapuram_2011}, healthy-ageing literature, frailty theory and research on social determinants can be integrated without collapsing their distinct insights. Biomedical Potential captures intrinsic physiological and functional resources; Social Modifiers capture the conditions through which those resources are supported or constrained; and Systemic Burdens capture structural barriers that limit the realisation of available health potential.

The conceptual formula introduced in that chapter is a theoretical statement, not a directly observed quantity. Its purpose is to formalise the proposition that future health capability depends jointly upon biomedical, social and systemic conditions, and the transition from that formula to the latent measurement model provided the bridge from theoretical argument to empirical evaluation. That bridge is set out formally below.

\begin{derivation}[From the conceptual formula to the latent measurement model]
\label{der:bridge}
The heuristic $H^{(c)}=(LQ)S-B$ asserts the direction and role of the three
domains; it is not the empirical scoring equation. Let
$\boldsymbol\eta_i=(\eta_{Bi},\eta_{Si},\eta_{Yi})^{\mathsf T}$ denote Biomedical
Potential, Social Modifiers, and the reverse-oriented Systemic domain for person
$i$, so larger values of every component indicate greater capability. Then
\[
 \mathbf y_i=\boldsymbol\tau+\boldsymbol\Lambda\boldsymbol\eta_i+
 \boldsymbol\varepsilon_i,
 \qquad
 \boldsymbol\eta_i=\boldsymbol\alpha+\boldsymbol\Gamma\zeta_i+
 \boldsymbol\delta_i .
\]
Fixing $\operatorname{Var}(\zeta)=1$ identifies the scale and gives
\[
 \operatorname{Cov}(\mathbf y)=\boldsymbol\Lambda
 (\boldsymbol\Gamma\boldsymbol\Gamma^{\mathsf T}+\boldsymbol\Psi_\delta)
 \boldsymbol\Lambda^{\mathsf T}+\boldsymbol\Theta.
\]
Theory determines the factor pattern, orientations, and expected signs;
$\boldsymbol\tau,\boldsymbol\Lambda,\boldsymbol\Gamma,
\boldsymbol\Theta,$ and $\boldsymbol\Psi_\delta$ are estimable. The predictor is
$\widehat H_i=E(\zeta_i\mid\mathbf y_i,\widehat{\boldsymbol\theta})$, not
$H^{(c)}$. Because only three first-order factors inform the higher-order model,
its global fit may be just identified under common constraints; correlated-factor
and bifactor \citep{reise_2012} alternatives must therefore be compared before accepting one score.
\end{derivation}

\begin{reading}[Why the conceptual formula is not the score]
$H^{(c)}=(LQ)S-B$ was written to say three things: that biology sets a ceiling,
that social conditions scale it, and that burdens subtract from it. It was
never meant to be evaluated for an individual. If it were, Profile~A would
need $L$ and $Q$ read off an instrument that does not exist, and the
subtraction of $B$ would let a large enough social score cancel the structural
burden outright.

What is computed instead is $\widehat H_i$, the model's best estimate of a
single unobserved quantity given a respondent's answers. Being a conditional
expectation it carries an error bar, and it is pulled towards the population
mean when the pattern of answers is unusual. A profile such as
$(0.55,0.35,0.40)$, Profile~A, comes out of that estimation, not out of the
heuristic.

The last sentence of the derivation is a warning about arithmetic. Three
first-order factors leave a higher-order model with no spare degrees of freedom
under the usual constraints, so its fit statistics can look perfect while
saying nothing at all. A correlated-factor model and a bifactor model have to
be fitted alongside it before anyone accepts one number standing in for three
domains.
\end{reading}

Aim~1 was therefore achieved at the conceptual level. The thesis produced a defined, internally coherent and empirically testable theory of health capability. The later empirical results determine how successfully the initial operationalisation represents that theory, but they do not remove the conceptual contribution itself.

A coherent theory is not yet a measurable one. The second aim is where the
theory was exposed to data.

\subsection{Aim 2: Develop and Validate the Measurement Model}

The second aim was to operationalise the framework as a multidimensional latent construct and evaluate its measurement properties. This measurement validation is the study's primary outcome. Chapter~\ref{ch:results1} addressed this aim through estimation of the proposed second-order structural equation model. The measurement analysis evaluated whether the observed indicators loaded onto the hypothesised first-order domains and whether those domains formed a coherent higher-order $H_{\text{potential}}$ construct, assessing model fit, reliability, convergent validity, discriminant validity and the stability of the estimated factor scores according to the pre-specified analysis plan.

\medskip
\noindent\emph{[Placeholder: state the final verdict on Aim~2. Report whether the three-domain structure achieved acceptable fit; whether Biomedical Potential, Social Modifiers and Systemic Burdens loaded meaningfully onto the higher-order construct; whether reliability and validity criteria were met; and whether any modifications were required. Distinguish clearly between full, partial and unsupported confirmation of the measurement hypothesis.]}

The answer to Aim~2 should be stated proportionately. If the hypothesised model was supported without major modification, the conclusion is that the $H_{\text{potential}}$ can be represented empirically as a higher-order latent construct in the NHATS population. If support was partial, the conclusion should identify which domains or indicators performed as expected and which require refinement. If the model was not supported, the conclusion should state that the proposed operationalisation did not adequately represent the theoretical construct in this cohort, while distinguishing failure of the measurement model from failure of the broader conceptual theory. This distinction is essential, because structural equation modelling tests a specific operational representation of the construct, not every possible way of measuring future health capability.

Each verdict has a different consequence for what follows, and the order of the
aims was chosen so that the dependence runs one way only. If the higher-order
structure holds, Aim~3 tests a single score whose meaning is settled. If the
three domains hold but the general factor is not determinate enough to score
individuals, a result the introduction anticipated, then the prognostic analysis
uses the domain scores. The realisation model is unaffected in that case, since
it takes the three domains as its inputs; what is withdrawn is the single
higher-order number. If the domains themselves fail, Aim~3 loses its object: its
results would be evidence about a set of survey items, not about health
capability.

Whether the construct can be measured and whether it is worth measuring are
different questions. The third aim asks the second of them.

\subsection{Aim 3: Evaluate Prognostic Performance}

The third aim was to test, as confirmatory evidence, whether the estimated $H_{\text{potential}}$ predicts clinically important future outcomes and adds value beyond established measures. Chapter~\ref{ch:results2} addressed this aim using five-year follow-up for two co-primary endpoints analysed separately, all-cause mortality and first overnight hospital stay, with a combined endpoint kept only as a sensitivity analysis. The prognostic analysis examined the association between the standardised latent score and time to event for each endpoint, first without adjustment and then after adjustment for age and sex; it also evaluated discrimination using Harrell's and Uno's concordance \citep{uno_2011} statistics, calibration through calibration-in-the-large and the calibration slope, comparative performance against established measures, and potential clinical utility through decision-curve analysis.

\medskip
\noindent\emph{[Placeholder: state the final verdict on Aim~3. Report whether a higher $H_{\text{potential}}$ was associated with lower risk of each co-primary outcome; whether the associations remained after adjustment; whether discrimination and calibration were acceptable; whether optimism correction materially reduced performance; and whether the construct was superior, non-inferior or inferior to the pre-specified comparators.]}

The conclusion should distinguish association from predictive usefulness. A statistically significant hazard ratio would show that the construct is related to future outcomes, but it would not by itself establish strong discrimination, accurate calibration or meaningful clinical utility. Conversely, modest incremental discrimination may still be important if the model is well calibrated and provides greater net benefit \citep{vickers_elkin_2006} across relevant decision thresholds. The answer to Aim~3 should therefore be based upon the complete performance profile rather than a single statistic.

The answer must also allow for one split verdict. The head-to-head comparison of
the realisation model against the additive composite and the factor score can go
against the model while the factor score itself predicts well. That outcome
would support the measurement and narrow the formal contribution. It would show
that the non-compensatory form does not improve prediction in this cohort. It
would not show the axioms false, because an axiom about how deficits combine is
a claim about realisation, and death or admission within five years is only one
of its consequences. The axioms would then stand as a normative proposal whose
predictive advantage has not been shown.

Predicting well in a cohort of older Americans says nothing about
haemophilia. The fourth aim was set at a lower level of ambition for that
reason.

\subsection{Aim 4: Assess Translational Relevance to Haemophilia}

The fourth aim was to assess whether the general framework could be translated coherently into haemophilia. Chapter~\ref{ch:translational} addressed this aim conceptually. The chapter showed that the three domains can be mapped onto clinically meaningful haemophilia indicators without redefining the construct. Biomedical Potential can be represented through joint health, bleeding phenotype, functional mobility, inhibitor status and age-related comorbidity. Social Modifiers can be represented through access to comprehensive care, caregiver and peer support, health literacy and self-management. Systemic Burdens can be represented through treatment availability, supply-chain stability, financial coverage and structural barriers to specialist care. The chapter also specified a modified Delphi consultation to establish face and content validity before empirical modelling, followed by a roadmap for cohort assembly, factor-model re-estimation, measurement-invariance testing, recalibration, prognostic evaluation and governance.

Aim~4 was therefore achieved at the level intended by the thesis. The framework was shown to be conceptually applicable to ageing people with haemophilia, and a defensible pathway for future empirical validation was defined. The aim did not seek to establish that the haemophilia adaptation is already psychometrically or prognostically valid; those questions belong to the next phase of the research programme.

The limits of this answer should be stated plainly. Conceptual applicability is
the easiest of the four verdicts to reach, because no data could have
contradicted it. What gives it content is that the roadmap names the tests that
could later contradict it: a Delphi panel that rejects part of the indicator
mapping (Derivation~\ref{der:consensus}), loadings that differ across groups of
people with haemophilia (Derivation~\ref{der:invariance}), or a recalibrated
model that still misjudges risk in the new population
(Derivation~\ref{der:recalibration}).

The aims describe what the work attempted. They do not map one to one onto the
questions Chapter~\ref{ch:introduction} posed, and a question can stay open
after the aim that addressed it has been met.

% =====================================================================
\section{The Research Questions, Answered as Far as the Thesis Allows}
\label{sec:conc-questions}
% =====================================================================

\noindent Chapter~\ref{ch:introduction} posed one primary question, five
secondary questions and six hypotheses, H1 to H6. This section says how much of
each answer the thesis already holds and what each possible result would mean.
Placeholders mark what waits on the analysis.

\subsection{The Primary Question}

The primary question asked whether a theoretically specified construct of
health-related capability can be measured reliably and validly in later life
without collapsing into an existing measure of health status, frailty, social
vulnerability or general wellbeing. The thesis can answer half of this now. It
has shown how the question would be settled: by the fit and determinacy of the
measurement model, by the reliability of the score, and by discriminant validity
against each of the named measures. It cannot yet say whether the answer is yes.

\medskip
\noindent\emph{[Placeholder: state the answer to the primary research question, giving the reliability of the score and its correlation with each pre-specified comparator.]}

The words ``without collapsing'' carry most of the weight. A construct can fit
well and still correlate so closely with a frailty index that it adds only a
new name. If that happens, the honest answer to the primary question is
no, whatever the fit statistics say: the three domains, as measured in NHATS,
would reduce to a measure already in use.

\subsection{The Secondary Questions}

The first secondary question, tied to H1, asked whether the domains form a
coherent but distinguishable structure. Coherence and distinctness pull in
opposite directions. Indicators that cohere strongly within a domain may also
correlate across domains, and if the domain correlations approach one the three
domains are one domain under three names. The test is passed when within-domain
coherence is adequate and the domain correlations stay clearly below the level at
which the domains cannot be told apart.

The second, tied to H2, asked whether a higher-order representation is
warranted. Fit alone cannot answer this (Derivation~\ref{der:bridge}), so the
answer rests on factor determinacy and on comparison with the correlated-factor
and bifactor models. If the general factor is not determinate, the answer is no,
and domain scores are reported rather than a total alone.

\medskip
\noindent\emph{[Placeholder: state the answers to the first and second secondary questions, identifying the measurement model retained and the evidence for retaining it.]}

The third secondary question, tied to H3, asked whether the measure adds
prognostic information. Its answer is the answer to Aim~3, and the conditions
stated there apply.

The fourth secondary question, tied to H4, asked whether the measurement
properties are stable across major demographic and socioeconomic groups. This is
the question on which fair use depends. If loadings or thresholds differ by
group, the same score means different things in different groups, and comparing
groups on it compares unlike quantities. A construct that includes social and
systemic domains is more exposed to this than a purely biomedical one, because
the meaning of an item about access or support may itself vary with social
position.

\medskip
\noindent\emph{[Placeholder: state whether measurement invariance was supported across the pre-specified groups, at which level, and which items or groups were non-invariant.]}

The fifth secondary question, tied to H5, asked whether domain profiles carry
information that a single score loses. The formal half of the answer is already
given. The exact decomposition of Chapter~\ref{chap:crf} separates any
capability gap into biomedical, social and structural parts, and the reading of
Profile~B in \S\ref{sec:future} shows a low structural coordinate pointing at the
payer rather than the clinic, which a total alone would not reveal. The empirical half, whether the profiles also predict outcomes
better than the total, waits on the analysis. If the profiles matter for action but add
nothing to prediction, they should be reported for diagnosis and the total used
for risk.

\medskip
\noindent\emph{[Placeholder: state whether domain-level profiles added information beyond the single score, for prediction and for identifying the source of constraint.]}

H6, on the head-to-head comparison of the Capability Realisation Model, is
answered with Aim~3 above. H4 and H5 have no aim of their own, which is why they
are answered here.

The questions ask what the thesis can show. The hypotheses state what it
expected to find, and they are answered more tersely because each admits fewer
intermediate verdicts.

% =====================================================================
\section{Status of the Research Hypotheses}
\label{sec:hypotheses}
% =====================================================================

\noindent The thesis was guided by hypotheses concerning the measurement structure, prognostic performance and translational relevance of the $H_{\text{potential}}$. They are grouped here by what each would establish. The first below corresponds to H1 and H2 of Chapter~\ref{ch:introduction}, the second and third to H3, and the fourth to the translational objective; H4, H5 and H6 are answered in \S\ref{sec:conc-questions}.

The first empirical hypothesis (H1 and H2), the study's primary hypothesis, concerned whether Biomedical Potential, Social Modifiers and Systemic Burdens form a valid multidimensional higher-order construct.

\medskip
\noindent\emph{[Placeholder: state whether the measurement hypothesis was supported, partially supported or not supported, and identify the decisive evidence.]}

The second empirical hypothesis (the association part of H3), the first of the confirmatory hypotheses, concerned whether the $H_{\text{potential}}$ predicts five-year all-cause mortality and, separately, first overnight hospital stay.

\medskip
\noindent\emph{[Placeholder: state whether the prognostic association hypothesis was supported and report the direction, magnitude and precision of the association.]}

The third empirical hypothesis (the incremental part of H3), also confirmatory, concerned whether the construct provides predictive information beyond conventional measures.

\medskip
\noindent\emph{[Placeholder: state whether the incremental-value hypothesis was supported, partially supported or not supported, drawing on discrimination, calibration, information criteria and decision-curve analysis.]}

The fourth hypothesis, which belongs to the translational objective rather than to H1 to H6, concerned whether the general framework is conceptually applicable to haemophilia. This hypothesis was supported at the conceptual and methodological level. The three domains can be mapped to clinically relevant haemophilia indicators, the mapping possesses a clear biological and health-system rationale, and the proposed Delphi and validation pathway provides a structured method for future empirical testing. This conclusion remains one of conceptual applicability, not completed clinical validation.

Verdicts of this kind are worth no more than the conditions attached to them.
Those conditions are set out next.

% =====================================================================
\section{Scope of the Conclusions}
\label{sec:scope}
% =====================================================================

\noindent The contribution of the thesis should be interpreted within the boundaries of the study design. Each limit below is a statement about where the boundary of the evidence lies, in the sense of \S\ref{sec:sys-boundary}: what was inside the system studied, and therefore what a conclusion can be about. The $H_{\text{potential}}$ was developed and internally validated in a single cohort of community-dwelling adults aged 65 years and older in the United States, so the empirical findings apply most directly to that population and setting. Bootstrap \citep{efron_tibshirani_1993,steyerberg_2019} validation provides an estimate of optimism and likely performance in similar samples, but it does not establish transportability to different countries, age groups, healthcare systems or disease populations; external validation remains necessary before the measurement and prognostic findings can be generalised more broadly.

The observational design also limits causal interpretation. The $H_{\text{potential}}$ is evaluated as a prognostic construct, and associations between the score and future outcomes do not demonstrate that interventions that alter the score will necessarily alter those outcomes. Establishing causality would require explicit causal designs, stronger assumptions and, ultimately, prospective intervention studies.

The use of estimated factor scores introduces additional uncertainty. Although the full measurement-and-prediction pipeline is repeated within the bootstrap procedure, the prognostic model still relies on an estimated rather than directly observed predictor. Future work may examine joint latent-survival models, Bayesian approaches or other methods that propagate measurement uncertainty more directly, along the lines set out below.

\begin{derivation}[Propagating factor-score uncertainty into the prognostic model]
\label{der:uncertainty}
Write $\widehat H_i=H_i+u_i$, with $E(u_i\mid H_i)=0$ and reliability
$\rho_H=\operatorname{Var}(H)/\operatorname{Var}(\widehat H)$. Under the
classical-error approximation, treating $\widehat H$ as observed produces
$\beta_{\mathrm{naive}}\approx\rho_H\beta$ and omits uncertainty from
estimating the measurement parameters and scores \citep{prentice_1982}. The
approximation is diagnostic rather than exact for Cox regression. A joint model combines
\[
 p(\mathbf y_i\mid H_i,\boldsymbol\theta_m),
\]
\[
 h_i(t\mid H_i,\mathbf z_i)=h_0(t)
 \exp\{\beta H_i+\boldsymbol\gamma^{\mathsf T}\mathbf z_i\},
\]
and integrates the likelihood over $H_i$. The primary two-stage analysis should
refit imputation, SEM, scoring, and survival modelling inside every bootstrap
replicate; a joint latent-survival or plausible-values model is the preferred
sensitivity analysis.
\end{derivation}

\begin{reading}[What attenuation costs, in numbers]
Reliability $\rho_H$ is the share of the variance in the estimated score that
is signal. Say it is $0.80$, which would be respectable for a factor score
built from a dozen indicators. If the true log hazard ratio per standard
deviation of $H$ is $-0.40$, a hazard ratio of $0.67$, then treating the
estimate as though it were the quantity itself returns about $-0.32$, a hazard
ratio of $0.73$. The association is real, its direction is right, and its
strength is understated by roughly a fifth.

This cuts two ways. A head-to-head comparison with a directly counted
competitor, a frailty index built from recorded deficits, is unfair to the
$H_{\text{potential}}$, since only one of the two scores is attenuated. It also
means the construct may not be credited with predictive strength it has not
shown: the report gives the attenuated figure, with the correction given as
a sensitivity analysis.

The joint model removes the problem by refusing to separate the two stages. It
writes the measurement and the survival model as one likelihood and integrates
over the unknown $H_i$, so uncertainty in each respondent's score carries into
uncertainty about the predicted risk. It costs a great deal more computation
and it lets errors in the measurement model contaminate the survival estimates,
which is why it is proposed here as a check on the two-stage analysis and not
as the analysis itself.
\end{reading}

The case that matters for the conclusions of this thesis is not the one above,
where the association survives. It is the one where it does not.

\begin{workedexample}[conc:attenuation]{How a null result should be read at a given reliability}
Suppose the analysis returns no association: a hazard ratio for the
$H_{\text{potential}}$ whose confidence interval includes $1$. Suppose also
that the reported reliability of the factor score is $\rho_H=0.55$, which would
be poor but not unheard of for a score built from few indicators per domain.
The question is what that null permits anyone to say.

Work backwards. If the true log hazard ratio per standard deviation is
$\beta=-0.22$, a hazard ratio of $0.80$, then
Proposition~\ref{prop:attenuation} gives
$\beta_{\mathrm{naive}}\approx0.55\times(-0.22)=-0.121$, a reported hazard ratio
of $0.886$. Take the standard error of the reported coefficient as $0.075$. The
reported interval is then $0.765$ to $1.026$, which includes $1$: a null. Had
the same true effect been measured without error, the interval would have run
$0.693$ to $0.930$ and excluded $1$.

The reliability that would have made the difference can be read off directly.
The interval clears $1$ when $\rho_H|\beta|>1.96\times0.075=0.147$, that is when
$\rho_H>0.147/0.22=0.67$. So a null at $\rho_H=0.55$ is consistent with a true
hazard ratio of $0.80$, while a null at $\rho_H=0.90$ would not be.

\begin{wecode}[Python]
import math
from hpv.prognostic import attenuated_coefficient

beta_true, se = -0.22, 0.075
for rho in (1.00, 0.55):
    b = attenuated_coefficient(beta_true, rho)
    lo, hi = math.exp(b - 1.96 * se), math.exp(b + 1.96 * se)
    print(f"rho {rho:.2f}: HR {math.exp(b):.3f}  95% CI {lo:.3f} to {hi:.3f}")

print(1.96 * se / abs(beta_true))   # reliability needed to clear 1: 0.67
\end{wecode}

\weanswer{At $\rho_H=0.55$ a true hazard ratio of $0.80$ is reported as $0.886$
with an interval of $0.765$ to $1.026$, so the null is uninformative about an
effect of that size. A reliability above about $0.67$ would have been needed to
rule it out. Asserted by
\texttt{test\_worked\_example\_conc\_null\_result\_at\_low\_reliability}.}
\end{workedexample}

\begin{reading}[The rule this fixes for the concluding chapter]
A null cannot be reported on its own. It has to be reported with the
reliability of the score and with the effect size the study could have detected
at that reliability, because those three numbers together are the finding and
any one of them alone is not.

The asymmetry is worth stating plainly. A positive result is understated by
attenuation and is therefore safe to believe as a lower bound. A null is
weakened by it, and the weaker the measurement the less a null means. The
uncomfortable consequence is that a poorly measured construct is hard to refute,
which is a reason to report $\rho_H$ prominently rather than a reason to take
comfort from a null.
\end{reading}

The available NHATS indicators also constrain the empirical form of the construct. The theoretical framework is broader than any one dataset, and some aspects of social or systemic capability may be represented only imperfectly in the developmental cohort; external studies with richer contextual, biological and health-system information may therefore refine the measurement model without abandoning the underlying theory. Finally, the framework contains unavoidable normative choices. Decisions about which domains belong within health, how disadvantage is coded and how biological, social and systemic dimensions should contribute to a common score cannot be derived from statistical evidence alone. The thesis makes these choices explicit and grounds them in established theory, but they remain open to ethical, philosophical and policy debate. These limitations do not negate the contribution; they define the stage the research has reached. The thesis establishes a theory, an initial operationalisation, an internal empirical test and a route to translation. It does not claim to deliver a universally valid clinical score or a final theory of health.

Naming a limit says what the evidence cannot support. It leaves open a practical
question: what may be done with the instrument now, before any limit has been
removed?

% =====================================================================
\section{Recommendations for Use}
\label{sec:conc-use}
% =====================================================================

\noindent The answer depends on what the use could do to the person scored. A
use that informs a research question and changes nothing for any individual can
tolerate an instrument tested in one cohort. A use that changes what a person
receives cannot. The recommendations below follow that line.

\subsection{Uses the Present Evidence Can Support}

Once the analysis is complete, and subject to its verdicts, the
$H_{\text{potential}}$ may be used as a research measure in cohort studies of
later life, where the question is how biological, social and structural
conditions jointly relate to later outcomes, and to generate hypotheses about
where constraint lies, by reading the domain profiles of groups with similar
totals. In such uses an error in the score costs a wrong lead, not a withheld
service.

\subsection{Uses That Must Wait}

The score should not be used to make or support decisions about individuals:
clinical triage, eligibility for a service, priority on a waiting list, or the
level of a benefit. Nor should population scores be used to allocate budgets
between regions or programmes. Three reasons combine. The score has been estimated in one cohort and has no
external calibration. Its normative choices, which domains count and how
disadvantage is coded, are open to dispute and have not been tested with the
people affected. And a score that includes structural burden can feed back on
the person scored, a risk taken up in \S\ref{sec:conc-systems}.

The haemophilia adaptation has a stricter status still. No haemophilia data have
been analysed, so there is no haemophilia score to use. Until the Delphi
consultation and a cohort study have been completed, the domain mapping of
Chapter~\ref{ch:translational} is a design for an instrument, not an instrument.

\subsection{Rules for Reporting the Score}

Whenever the score is used, including in research, the following should be
reported with it.

\begin{enumerate}
  \item The three domain scores alongside any total, so that a reader can see
  which domain drives a low value.
  \item The reliability of the score in the sample used and, with any null
  result, the smallest effect that could have been detected at that
  reliability, as the worked example in \S\ref{sec:scope} requires.
  \item Calibration in the population of use as well as discrimination, since
  a score can rank people correctly while misstating their risk.
  \item The version of the measurement model and the frozen values of $\alpha$,
  $\rho$ and the normalisation bounds, so that results from different versions
  are not pooled.
  \item The inequality-aversion parameter, whenever the equity-sensitive
  population functional of Chapter~\ref{chap:crf} is used, since the ranking of
  policies can change with it.
\end{enumerate}

\subsection{Conditions for Moving from Research to Practice}

Practical use would require, at a minimum, external validation and
recalibration in the population where the score would be used; evidence of
measurement invariance across the groups between which decisions would be made;
a named body responsible for the model, with the monitoring duties set out in
\S\ref{sec:conc-systems}; a route by which a person can see and challenge the
domain breakdown of their score; and a rule, stated in advance, for suspending
use when monitoring shows the model has drifted. None of these is a statistical
result.

These conditions say what must be in place. What would put them there is
further research, taken next.

% =====================================================================
\section{Implications for Future Research}
\label{sec:future}

\subsection{Economic and policy limits}

The framework's policy value lies first in diagnosis, not allocation. Domain
profiles can distinguish a clinical-capacity deficit from inadequate support or
an institutional barrier and can therefore help identify which sector bears
responsibility for a response. A scalar score cannot make that distinction once
the profile is discarded.

Four questions must remain separate. Need asks who is constrained. Effectiveness
asks whether an intervention changes a relevant outcome. Cost-effectiveness asks
whether its gains exceed the health opportunity cost of its resources.
Affordability asks whether the required expenditure can be accommodated.
H-Potential principally informs the first and may support heterogeneity analysis
for the second. It does not answer the latter two.

\begin{reading}[The four questions, asked about Profile~B]
Need. With structure unfavourable the profile is $(0.30,0.70,0.20)$. The low
third coordinate, not the low first, is what the framework flags, and it points
at the funding of treatment. The body that can act is the payer, not the
haemophilia centre. This is the diagnostic step, and it is the part the
$H_{\text{potential}}$ is built to do.

Effectiveness. Would better access change bleeding, joint health or
admissions? The framework does not answer that. A trial or a carefully
designed observational study does.

Cost-effectiveness. Money spent on the treatment buys health that some other
patient in the same budget then does not get. Whether the exchange is worth
making turns on the health opportunity cost in that system, about which the
profile says nothing.

Affordability. A favourable ratio can still exceed what a provincial budget
carries in a given year. That is a question about cash and timing.

Running these four together is how a measure of unmet need turns into a
rationing rule. Profile~B with structure unfavourable scores low in part
because the health system has failed to supply treatment. A score read without
the profile behind it could be taken as evidence of a poor investment, which
inverts the finding.
\end{reading}

Future economic work should therefore embed profile change within conventional
and distributional cost-effectiveness analysis rather than define a
cost-per-H-Potential threshold. It should estimate intervention effects and
costs by relevant groups, include the distribution of displaced health and
state openly any inequality-aversion parameter. \citep{asaria_2016_dcea} This is
consistent with WHO's global account of HTA as a multidisciplinary process, but
the resulting decision cannot simply be exported from one health system to
another. \citep{who_hta_2014,who_hta_institutional_2021} Prices, service
capacity, the marginal productivity of expenditure, fiscal limits, legal
entitlements and distributive priorities differ by jurisdiction. Even when an
intervention produces the same domain change, its opportunity cost and policy
ranking may therefore differ.

The United States is the empirical development setting, not the normative
centre of the framework. Application elsewhere requires conceptual review,
local item adaptation, invariance testing, external prognostic validation and
jurisdiction-specific economic evaluation. South Africa is one valuable future
test because its distribution of resources and emerging HTA institutions may
make systemic constraints especially visible; it is an illustration of
translation, not the geographical boundary of the theory.
\citep{south_africa_hta_2026}

Policy use also requires governance beyond statistical validation. Public
specification, domain-level explanation, community participation, periodic
review, disparate-impact monitoring and a meaningful appeal route are parts of
the validity of the practice, not administrative additions to it. A low score
caused by structural constraint is evidence of a potential institutional duty
to act, never evidence that the person is less deserving of investment.

Governance of that kind is a long-term requirement. The immediate work is
narrower and empirical.
% =====================================================================

\noindent The immediate next step is external validation of the measurement and prognostic models in independent cohorts. Such validation should assess the stability of the three-domain structure, the reproducibility of factor loadings, the calibration of predicted risks and the need for population-specific recalibration. Validation across health systems is especially important because Systemic Burdens are expected to vary substantially by institutional and national context. A model developed in the United States may require material adaptation in settings characterised by different financing structures, levels of social protection or access to healthcare; these differences are not merely sources of statistical heterogeneity but are central to the theory being tested.

The proposed haemophilia programme represents the first disease-specific extension. The next phase should conduct the modified Delphi consultation, assemble an appropriate cohort and evaluate whether the translated indicators reproduce the intended domains. Only after satisfactory psychometric estimation should the haemophilia-specific score be tested against bleeding, joint, hospitalisation, quality-of-life and mortality outcomes. Further disease-specific adaptations could examine whether the same architecture remains useful in other chronic conditions, such as sickle cell disease, chronic kidney disease, diabetes, neurodegenerative disorders and multimorbidity, each of which involves distinct biological processes but also strong social and systemic influences.

A longitudinal extension is also required. The present work estimates $H_{\text{potential}}$ at a defined baseline point, but future capability is unlikely to remain static. Repeated measurement could distinguish stable between-person differences from within-person change, identify trajectories of decline or recovery, and evaluate whether changes in the latent construct precede clinical events. The dynamical architectures described in Appendices~\ref{app:dynamical} and~\ref{app:hemo_dynamics} provide possible directions for modelling these processes: ordinary, stochastic and partial differential equations may eventually support representations of health potential that evolve continuously under biological, social and treatment inputs. Those models remain theoretical extensions and should be developed and validated separately from the present measurement-and-prediction programme.

Repeated measurement would also allow a test that the static design cannot make.
Work on critical transitions proposes that a system close to a threshold
recovers more slowly from small disturbances \citep{scheffer_2009}, and slowing
of recovery has been proposed as a general marker of risk of acute transitions
in chronic disease \citep{olde_rikkert_2016}. Physical resilience in older
adults has likewise been framed as a construct in its own right, defined by the
response to a stressor rather than by the level of function before it
\citep{whitson_2016}. If health capability behaves as a reserve in the sense of
\S\ref{sec:sys-resilience}, the speed with which domain scores return after a
hospital stay or a loss of support may carry information that the level of the
score does not. A cohort measured before and at intervals after such events
could test whether recovery speed adds to the prediction made from level alone.

The order of these tasks matters. External validation comes first because every
later task inherits its answer. A dynamical model of a quantity that has not
been shown to transport carries the errors of the static model into every time
step. Economic and policy work comes last, because it needs the effects of
interventions on domains that are known to be measured consistently across the
groups compared. Each task also has a stated way to fail: external
validation fails if recalibration does not restore calibration; the haemophilia
work fails if the Delphi panel or the invariance tests reject the mapping; the
longitudinal work fails if change in the score does not precede change in
outcomes. Stating these in advance keeps each step a test rather than an
elaboration.

Future research should also address practical interpretability. Latent scores are mathematically useful but may be difficult to communicate to clinicians, patients and policymakers, and work will be needed to establish meaningful reference ranges, identify interpretable risk strata and determine whether domain-specific profiles are more useful than a single overall score in some settings. Ethical and distributive implications require particular attention. Because social disadvantage and systemic barriers contribute to the construct, use of the score in resource allocation could either reveal inequity or inadvertently reproduce it. A low score arising from poor healthcare access should not be interpreted as an intrinsic deficiency of the individual, and any practical use must therefore distinguish between identifying unmet need and penalising those already disadvantaged.

These are tasks the framework sets for others as much as for its author. They
are framed as research. One question they leave open is not a research question
at all: if the score were deployed, who would watch it, and how would they know
it had gone wrong?

% =====================================================================
\section{A Systems View of the Conclusions}
\label{sec:conc-systems}
% =====================================================================

\noindent The sections above state the conclusions as claims about a construct
and a score. Chapter~\ref{ch:systems} argued that the $H_{\text{potential}}$
sits inside two systems at once: the system of biological, social and structural
conditions whose capacity it describes, and the measurement pipeline that
produces the number. This section reads the conclusions through that view. It
asks first what the systems reading adds to them, and then what monitoring and
governance a deployed score would need.

\subsection{What the Systems Reading Adds}

The first addition concerns the boundary. Validity is not a property a score
carries with it; it holds for a stated population, data source and health
system, which together form the boundary of the system studied
(\S\ref{sec:sys-boundary}). Read this way, the limits in \S\ref{sec:scope} stop
being a list of cautions and become a prediction. The Systemic Burdens domain is
the part of the construct most tied to its environment, since its indicators
describe financing, access and services that differ between health systems. When
the instrument is moved across a boundary, that domain should be the first to
lose invariance or calibration, and external validation can test this directly.

The second addition is the line between verification and validation
(\S\ref{sec:sys-pipeline}). Verification asks whether the pipeline was built as
specified; validation asks whether it measures and predicts what it was meant
to \citep{incose_2015}. The conclusions of this chapter divide along that line.
Every statement made here without a placeholder is either conceptual or a matter
of verification: the axioms and their consequences, the arithmetic of the worked
examples, the logic of the design. Appendix~\ref{app:verification} checks the
computable part of it by code. Every statement with a placeholder is a matter of
validation, and no amount of verification can stand in for it.
The third addition concerns defence in depth and its limits
(\S\ref{sec:sys-fmea}). The design places several layers between the data and a
false conclusion: a construct defined without the outcomes, parameters frozen
before the outcomes are seen, every stage refitted in each bootstrap replicate,
and comparison against pre-specified competitors. Layered defences work when the
layers fail for different reasons \citep{reason_1990}. These layers do not fully
meet that condition. All draw on one cohort and the same survey items. A defect
in the source data, such as the reversed response codes recorded in the data checks, passes through every layer at once, because every layer
trusts the same input. That is a common-cause failure. Appendix~\ref{app:datachecks}
records the checks that guard against the known cases, but no internal check
can guard against defects shared by the whole cohort. Only data from elsewhere
can. This gives external validation a second justification beyond
transportability: it is the one step that removes a whole class of failure the
internal checks cannot reach.

The fourth addition concerns feedback (\S\ref{sec:sys-loops}). The
complementarity theorem of the realisation model says that the return to
biomedical potential depends on social and structural conditions. In a system
that changes over time, that is the static trace of loops between the domains:
poor access may wear down biomedical reserve, lower reserve may raise the
demands on support, and strained support may in turn reduce access. A single cross-sectional score cannot
tell a reinforcing loop of this kind from a fixed difference between people.
Claims about loops must wait for the longitudinal work of \S\ref{sec:future}.

\subsection{Monitoring and Governance of a Deployed Score}

A score in use becomes part of a control structure (\S\ref{sec:sys-control}).
Controllers act on the system on the basis of a model of it, and they need
feedback to learn when their model is wrong \citep{leveson_2011}. For the
$H_{\text{potential}}$ the controllers would include the researchers who
maintain the measurement model, clinicians or programme staff who read
individual scores, payers and planners who read population scores, and the
people scored, who act on their own health and may contest how they are
described.

The main hazard follows from the construct itself. The score includes structural
burden, and the people acting on it can change that burden. If a low systemic
score leads a payer to supply treatment, the burden falls and the score rises.
The loop is balancing, and it is what the framework intends. If instead a low
total is read as a poor prospect and leads to less investment, the burden rises
and the score falls further. That loop is reinforcing, and it conceals its own
cause, since each round of decline appears to confirm the judgement that started
it. The governance requirement that follows is concrete: a low structural score
must be routed to the body that controls the structure, as the reading of
Profile~B in \S\ref{sec:future} showed, and must never be used to reduce what
the person receives.

A deployed score would need at least the following monitoring.

\begin{enumerate}
  \item \textbf{Calibration drift.} Predicted and observed event rates compared
  at set intervals in the population of use, with a threshold for
  recalibration fixed before deployment.
  \item \textbf{Measurement drift.} Re-estimation of loadings and repeated
  invariance checks whenever indicators, survey forms or coding change, since
  a changed item changes the construct without any change in the people.
  \item \textbf{Disparate impact.} Score distributions, and the actions taken
  on them, compared across the groups for which invariance was tested, so that
  harm can be seen in the actions as well as in the scores.
  \item \textbf{Action feedback.} A record of what was done in response to each
  class of score and what followed. This is the feedback path most often
  missing, and without it a controller cannot learn whether its actions help.
  \item \textbf{Appeal and correction.} A route by which a person can see the
  domain breakdown of their score and challenge it, with corrections fed back
  into the data checks.
  \item \textbf{Versioning and withdrawal.} Every score labelled with the model
  version that produced it, and a rule, stated in advance, for suspending use
  when calibration or invariance fails.
\end{enumerate}

\noindent Delay makes this harder than it looks (\S\ref{sec:sys-stocks}). The
outcomes the score predicts arrive years after it is computed; the endpoints in
this thesis run to five years. Calibration feedback is therefore slow, and a
controller may act for years on a model whose miscalibration it cannot yet see.
Measurement drift, data checks and appeals give faster signals than outcomes do,
which is why they belong in the monitoring plan as early warnings and not as
afterthoughts. Responsibility should also be divided clearly. The body that
maintains the model owns calibration, measurement drift and versioning; the body
that uses the score owns action feedback, disparate impact and appeals. A duty
that belongs to both tends to be done by neither.

\subsection{What the Systems Reading Cannot Add}

The systems reading adds requirements, not evidence (\S\ref{sec:sys-limits}). A
control structure drawn on paper does not show that anyone will act on its
signals, and a list of failure modes is only as complete as the imagination of
those who wrote it. All
models are wrong in some respect, and the useful question is which errors matter
for the use at hand \citep{box_1976}. The systems view helps to ask that
question. It does not answer it; the data and the people affected do.

What may be claimed before any of this is done is stated below.

% =====================================================================
\section{Final Concluding Statement}
\label{sec:final}
% =====================================================================

\noindent Subject to completion of the empirical analyses, this thesis offers a theory-grounded and empirically testable framework for measuring health as future capability. Its central contribution is the proposition that future health cannot be represented adequately through current disease status alone; it emerges from the interaction of biological reserve, social resources and the systems through which care is made available.

The $H_{\text{potential}}$ operationalises this proposition as a multidimensional latent construct. By estimating the measurement model before testing prognostic performance, the thesis preserves the distinction between what the construct means and what it predicts. By subsequently translating the domain structure to haemophilia, it shows how a general theory of health capability may be adapted to a specific disease without becoming reducible to that disease. At the present stage, the $H_{\text{potential}}$ remains a research instrument. It has been developed within a single cohort and requires independent external validation, recalibration and prospective evaluation before clinical or policy use can be justified, and the haemophilia application similarly remains a conceptual and methodological roadmap rather than an empirically validated disease-specific tool.

The wider contribution is therefore not a claim that one final metric has solved the problem of health measurement. It is the provision of a disciplined template for approaching that problem, summarised in Figure~\ref{fig:template}: begin with an explicit theory of health; distinguish latent constructs from their observed indicators; separate measurement from prediction; evaluate association, discrimination, calibration and utility; make normative choices visible; and require external validation before implementation.

\begin{figure}[htbp]
\centering
\begin{tikzpicture}[node distance=6mm and 6mm]
\tikzset{hpstep/.style={hpbox, text width=4.1cm, minimum height=1.45cm}}
\node[hpstep] (s1) {\textbf{1}\ Explicit theory of health};
\node[hpstep, right=of s1] (s2) {\textbf{2}\ Latent construct distinguished from its indicators};
\node[hpstep, right=of s2] (s3) {\textbf{3}\ Measurement separated from prediction};
\node[hpstep, below=of s3] (s4) {\textbf{4}\ Evaluation: association, discrimination, calibration, utility};
\node[hpstep, left=of s4] (s5) {\textbf{5}\ Normative choices made visible};
\node[hpstep, left=of s5] (s6) {\textbf{6}\ External validation before implementation};
\draw[hparrow] (s1) -- (s2);
\draw[hparrow] (s2) -- (s3);
\draw[hparrow] (s3) -- (s4);
\draw[hparrow] (s4) -- (s5);
\draw[hparrow] (s5) -- (s6);
\end{tikzpicture}

\captionsetup{justification=justified, width=\linewidth}
\caption{A reusable template for building theory-grounded, multidimensional health measures.}
\label{fig:template}
\end{figure}

If the empirical hypotheses are supported, the thesis will provide initial evidence that this template can produce a coherent and prognostically useful measure of future health capability. If parts of the empirical model require revision, the same framework will identify where the theory, indicators or estimation strategy must be refined. In either case, the work advances the argument that health measurement should attend to the capabilities, opportunities and constraints that shape what health can become, and not to realised burden alone.

What this chapter settles is the reading of every result the analysis could
return. It fixes, before the numbers are known, what a full, partial or failed
measurement model would mean; how a null must be reported against the
reliability of the score; what a split verdict between the factor score and the
realisation model would leave standing; what the instrument may be used for in
the meantime; and what must surround it before anyone acts on it. What it leaves
open is which of those readings applies, and that is answered by the analysis of
Chapters~\ref{ch:results1} and~\ref{ch:results2}, whose placeholders this
chapter shares. No chapter follows in the main text. The appendices hold the
mathematics, the dynamical extension, the code that verifies the computable
claims and the checks on the data, and the next piece of work lies outside the
thesis: running the analysis that fills the placeholders, and then testing the
instrument in a population it has not seen.
